\documentclass[11pt,aspectratio=169]{beamer}

\usetheme{Madrid}
\usecolortheme{seahorse}

\usepackage[utf8]{inputenc}
\usepackage{amsmath,amssymb}
\usepackage{algpseudocode}
\usepackage{algorithmicx}
\usepackage{xcolor}
\usepackage{listings}
\usepackage{booktabs}
\usepackage{tikz}
\usetikzlibrary{arrows.meta,positioning,fit,backgrounds}

\setbeamertemplate{navigation symbols}{}
\setbeamertemplate{footline}[frame number]

\title{Beyond the Syllabus}
\subtitle{Network Flow: The 66-Year Race Toward Almost-Linear Time\\
\small Week 10 Supplement -- TOAE209/TOAH209: Design and Analysis of Algorithms}
\author{Foreign Trade University}
\date{}

\begin{document}

% ---------------------------------------------------------------
\begin{frame}
  \titlepage
\end{frame}
% ---------------------------------------------------------------

\begin{frame}{Outline}
  \tableofcontents
\end{frame}

% =================================================================
\section{Motivation}
% =================================================================

\begin{frame}{Why look beyond the syllabus?}
  \begin{itemize}
    \item This week you proved that Edmonds-Karp -- Ford-Fulkerson augmenting along
    \emph{shortest} paths -- runs in $O(VE)$ augmentations, each costing $O(E)$ for the
    BFS, for a total of $O(VE^2)$: polynomial, and independent of the edge capacities.
    \item A natural question the lecture doesn't answer: \textbf{is $O(VE^2)$ close to
    the best possible, or is there room for a much faster algorithm?}
    \item The honest answer is startling: as of this decade, max-flow can be solved in
    time that is \textbf{almost linear} in the size of the input -- a bound that seemed
    out of reach for the 60+ years between Ford-Fulkerson (1956) and 2022.
  \end{itemize}
  \begin{alertblock}{How this supplement is different from a survey}
    We trace the \emph{actual sequence of ideas} -- not just headline numbers -- from
    this week's $O(VE^2)$ up to $m^{1+o(1)}$, with every specific bound and year sourced
    to a real, checkable paper.
  \end{alertblock}
\end{frame}

\begin{frame}{Why $O(VE^2)$ actually hurts at scale}
  \small
  Plug in graph sizes you'd meet in practice (web-scale routing, large logistics
  networks) and compare $O(VE^2)$ against ``almost linear'' ($m^{1+\epsilon}$ for tiny
  $\epsilon$, where $m \approx V \approx E$ here for simplicity):
  \begin{center}
  \begin{tabular}{@{}rrr@{}}
    \toprule
    $V = E = m$ & $VE^2 = m^3$ & $m^{1.01}$ (illustrative ``almost linear'') \\
    \midrule
    $10^3$ & $10^9$ & $\approx 1.1\times10^3$ \\
    $10^4$ & $10^{12}$ & $\approx 1.1\times10^4$ \\
    $10^6$ & $10^{18}$ & $\approx 1.4\times10^6$ \\
    \bottomrule
  \end{tabular}
  \end{center}
  \begin{exampleblock}{The gap}
    At $m=10^6$, $O(VE^2)=10^{18}$ operations is not just slow -- it is
    \emph{physically uncomputable} in a human lifetime on any real machine, while
    $m^{1.01}$ is finished before you can blink. This is the gap the algorithms below
    close.
  \end{exampleblock}
\end{frame}

% =================================================================
\section{A 70-Year Timeline}
% =================================================================

\begin{frame}{From Ford-Fulkerson to almost-linear time}
  \small
  \begin{center}
  \begin{tabular}{@{}lll@{}}
    \toprule
    Year & Algorithm & Bound \\
    \midrule
    1956 & Ford--Fulkerson & finite, but not polynomial in $n,m$ alone \\
    1972 & \textbf{Edmonds--Karp (this week)} & $O(VE^2)$ \\
    1998 & Goldberg--Rao & $O\big(m\min(\sqrt{m}, n^{2/3})\log\tfrac{n^2}{m}\log U\big)$ \\
    2022 & Chen, Kyng, Liu, Peng, Probst Gutenberg, Sachdeva & $m^{1+o(1)}$ (randomized) \\
    2023 & van den Brand \emph{et al.} & $m^{1+o(1)}$ (\textbf{deterministic}) \\
    2025 & van den Brand, Gholizadeh, Jiang, de Vos & parallel, $\tilde O(m{+}n^{1.5})$ work \\
    \bottomrule
  \end{tabular}
  \end{center}
  For 24 years (1998--2022), Goldberg--Rao's J.ACM bound was simply the best anyone
  could prove for general max-flow. Then two papers in three years broke it open.
\end{frame}

% =================================================================
\section{2022: Max-Flow in Almost-Linear Time}
% =================================================================

\begin{frame}{The 2022 breakthrough}
  \begin{block}{Chen, Kyng, Liu, Peng, Probst Gutenberg, Sachdeva -- ``Maximum Flow and
  Minimum-Cost Flow in Almost-Linear Time,'' FOCS 2022 (\textbf{Best Paper})}
    An algorithm computing \emph{exact} max-flow and min-cost flow on directed graphs
    with $m$ edges and polynomially-bounded integer capacities/costs, in
    $m^{1+o(1)}$ time.
  \end{block}
  \begin{itemize}
    \item $m^{1+o(1)}$ means: for any $\epsilon>0$, the running time is
    $O(m^{1+\epsilon})$ for $m$ large enough -- ``almost'' linear, with the hidden
    exponent shrinking to $0$ as $m$ grows. This is qualitatively different from any
    polynomial bound with a fixed exponent $>1$, like this week's $O(VE^2)$.
    \item arXiv:2203.00671; published at FOCS 2022, pp.\ 612--623.
  \end{itemize}
\end{frame}

\begin{frame}{The core idea: replace augmenting paths with a continuous method}
  \begin{itemize}
    \item Augmenting-path methods (Ford-Fulkerson, Edmonds-Karp) improve the flow
    \emph{one path at a time} -- inherently $\Omega(m)$ separate improvement steps in
    the worst case, each needing its own graph search.
    \item The 2022 algorithm instead builds the flow via an \textbf{interior-point
    method}: it takes only $m^{1/2+o(1)}$ \emph{global} steps, and each step is a
    single ``approximate minimum-ratio cycle'' computed on the \emph{whole} graph at
    once via an electrical-flow (graph Laplacian) solver.
    \item The remaining trick is a new \textbf{dynamic graph data structure} that
    maintains these approximate min-ratio cycles under edge-weight updates in
    \emph{amortized $m^{o(1)}$ time per update} -- structurally the same flavor of
    result as this week's ``near-constant amortized cost per operation,'' just for a
    much harder maintained quantity than Union-Find's connected components.
  \end{itemize}
\end{frame}

\begin{frame}{Why this result is bigger than just max-flow}
  \small
  The same interior-point-plus-dynamic-Laplacian-solver framework gives
  $m^{1+o(1)}$-time algorithms for any flow problem minimizing a \emph{general
  edge-separable convex cost} to high accuracy -- not just linear costs. The paper's
  own examples:
  \begin{itemize}
    \item entropy-regularized optimal transport (a core primitive in modern generative
    modeling and domain adaptation),
    \item matrix scaling,
    \item $p$-norm flows and $p$-norm isotonic regression on arbitrary DAGs.
  \end{itemize}
  \begin{alertblock}{The pattern to notice}
    A single algorithmic idea (fast dynamic Laplacian solvers) closed the gap on a
    32-year-old flow bound \emph{and} gave new best-known algorithms for four unrelated-
    looking problems at once -- exactly how the biggest wins in algorithms tend to
    happen.
  \end{alertblock}
\end{frame}

% =================================================================
\section{2023: Making It Deterministic}
% =================================================================

\begin{frame}{Why ``randomized'' was not the end of the story}
  \begin{itemize}
    \item The 2022 algorithm is \textbf{randomized}: correct with high probability, but
    it uses random bits and its worst-case guarantee is over the algorithm's own coin
    flips, not adversarial in the traditional deterministic sense.
    \item For safety-critical or adversarial settings, a \textbf{deterministic}
    guarantee -- correct on \emph{every} input, every time, with no dependence on luck
    -- is strictly stronger and often required.
  \end{itemize}
  \begin{block}{van den Brand, Chen, Kyng, Liu, Peng, Probst Gutenberg, Sachdeva, Sidford
  -- ``A Deterministic Almost-Linear Time Algorithm for Minimum-Cost Flow,'' FOCS 2023}
    Derandomizes the 2022 result: exact max-flow and min-cost flow in $m^{1+o(1)}$ time,
    \textbf{deterministically}. This was the first improvement to the deterministic
    max-flow running time (for polynomially-bounded capacities) since Goldberg--Rao in
    1998 -- a 25-year gap. arXiv:2309.16629.
  \end{block}
\end{frame}

% =================================================================
\section{2025: Making It Parallel}
% =================================================================

\begin{frame}{Why ``almost-linear time'' still isn't the end of the story}
  \begin{itemize}
    \item $m^{1+o(1)}$ counts \textbf{total work} on a single processor. On modern
    hardware (multi-core CPUs, GPUs, distributed clusters), the number that determines
    \emph{wall-clock} time is \textbf{depth} -- the length of the longest chain of
    operations that must happen strictly in sequence.
    \item An algorithm can have optimal total work and still be slow in practice if its
    depth is $\Omega(m)$ -- exactly the failure mode of augmenting-path methods, which
    are inherently sequential (each augmentation depends on the residual graph left by
    the last one).
  \end{itemize}
  \begin{block}{van den Brand, Gholizadeh, Jiang, de Vos -- ``Parallel Minimum Cost Flow
  in Near-Linear Work and Square Root Depth for Dense Instances,'' SPAA 2025}
    A randomized parallel min-cost-flow algorithm with $\tilde O(m+n^{1.5})$ total work
    and $\tilde O(\sqrt{n})$ depth, for graphs with $m > n^{1.5}$ edges -- the first
    algorithm combining near-linear work \emph{and} sublinear depth in that regime.
    Corollaries improve parallel bipartite matching, shortest paths, and reachability
    too. arXiv:2503.13274.
  \end{block}
\end{frame}

\begin{frame}{As of Aug 2026: still a very active frontier}
  \small
  This line of work has not slowed down. Titles accepted at ICALP 2026 continuing the
  same program include \emph{``Back in the Saddle: Toward Parallel Approximate
  Minimum-Cost Flow''} and \emph{``Parallel Reachability and Shortest Paths on
  Non-Sparse Digraphs: Near-Linear Work and Sub-Square-Root Depth''} -- both pushing
  the 2025 parallel result further. The exact final bounds for \emph{general} (sparse)
  graphs in \emph{low} parallel depth remain open as of this writing; treat any specific
  number you see quoted for 2026 results as something to verify against the published
  version, the same discipline this supplement has followed throughout.
\end{frame}

% =================================================================
\section{Summary}
% =================================================================

\begin{frame}{Summary}
  \footnotesize
  \begin{enumerate}
    \item \textbf{This week's $O(VE^2)$} is a correct, teachable, historically accurate
    stepping stone -- not the state of the art, and not intended to be.
    \item \textbf{The 24-year plateau (1998--2022)} at Goldberg--Rao's bound shows that
    ``no one has improved this in decades'' is not the same claim as ``no improvement
    exists.''
    \item \textbf{Almost-linear (2022) $\to$ deterministic (2023) $\to$ parallel
    (2025)}: three separate real results were needed to turn one breakthrough into a
    worst-case-safe, hardware-friendly algorithm -- finding the bound and making it
    fully deployable are different milestones.
    \item The engine (fast dynamic graph-Laplacian solvers) is the same kind of tool
    this course meets again informally in Week 15, and is why the field moved past
    purely combinatorial augmenting-path methods for the hardest instances.
  \end{enumerate}
  \begin{alertblock}{Looking ahead}
    Next week reduces \textbf{bipartite matching} to a single max-flow computation --
    in practice almost always solved with the simpler Hopcroft--Karp or Edmonds-Karp
    algorithms from this week, since real matching instances rarely need
    almost-linear-time machinery to run fast enough.
  \end{alertblock}
\end{frame}

\end{document}
